El desarrollo se organizó en iteraciones cortas y cerradas: cada una parte de una tarea
concreta, termina con el código, sus pruebas y la verificación automática superadas, y se
consolida en un \textit{commit} con su entrada en el registro de iteraciones
(\cref{AP:ITERATIONS}). El trabajo de programación lo ejecutaron dos \dfnpl{agente}, Claude
Code y Codex; las decisiones de alcance, de producto y de aceptación correspondieron al autor.
Esta sección describe las fases del proceso y deja explícito ese reparto, que se declara
también en \cref{SEC:IMP-AI} y en \cref{AP:AI}.

\subsection{Fases}

\begin{enumerate}
  \item \textbf{Prototipo previo.} Entre marzo y mayo de 2026 se construyó un primer
        prototipo contra un contrato preliminar del \dfn{backend} (rama
        \texttt{modulo-profesor}). Su contraste con el \dfn{backend} real motivó la
        reescritura descrita en \cref{SEC:INTRO-CONTEXT}; se conserva en una etiqueta de
        Git y no forma parte del resultado.
  \item \textbf{Análisis del contrato.} Se extrajo del documento OpenAPI real el
        inventario de operaciones y esquemas mediante \textit{scripts} deterministas, y
        se contrastaron las vistas de Django en los casos dudosos (\cref{SEC:AN-BACKEND}).
        El inventario resultante (\texttt{docs/frontend/API\_INVENTORY.md}) es desde
        entonces la única fuente admitida para las llamadas a la \acs{api}.
  \item \textbf{Base común.} Antes de repartir trabajo se construyó, en una única línea, la
        base compartida: sistema de diseño, cliente \acs{http}, simulación del
        \dfn{backend}, registro de módulos, inicio de sesión y los documentos de
        coordinación (\textit{commit} \texttt{1ebeac9}, 17 de octubre de 2026). Sin ella,
        el segundo agente no podía empezar sin modificar ficheros que el primero también
        necesitaba.
  \item \textbf{Desarrollo en paralelo.} A partir de la base se abrieron dos líneas sobre
        \dfnpl{worktree} distintos: el núcleo docente (corrección, exámenes, incidencias de
        escaneo, resultados y revisiones) y la gestión (administración de la organización y
        de asignaturas, perfil y notificaciones). El reparto se hizo por áreas funcionales
        con dependencias mínimas, no por volumen (\cref{SEC:IMP-PARALLEL}).
  \item \textbf{Iteraciones autónomas.} En la última fase, cada agente recibía un encargo
        priorizado y lo ejecutaba en bloques: tomar la siguiente tarea pendiente,
        completarla con pruebas, superar las comprobaciones, hacer el \textit{commit},
        actualizar su registro de progreso y terminar. El autor seguía el avance de forma
        remota y no intervenía en cada bloque; las dudas se resolvían con un criterio
        documentado y quedaban registradas para su revisión posterior.
\end{enumerate}

\subsection{Condición de cierre de una iteración}

Ninguna iteración se daba por terminada sin superar la misma secuencia de comprobaciones:
análisis estático (ESLint), pruebas (Vitest), compilación de producción (Vite) y auditoría de
estilo, agrupadas en la orden \texttt{npm run check}, seguidas de la auditoría de propiedad de
ficheros de la línea correspondiente (\cref{SEC:TEST-AUDITS}). La regla se endureció tras un
error de proceso: en una iteración temprana se encadenó la verificación con el
\textit{commit} sin comprobar el código de salida y se consolidó un estado con pruebas en
rojo. Se corrigió en el \textit{commit} siguiente y, a partir de entonces, el \textit{commit}
quedó condicionado al resultado de la verificación (\cref{AP:ITERATIONS}, entrada 13).

Las pruebas no se escribían al final: cada funcionalidad incorporaba sus casos en la misma
iteración, y varios defectos se descubrieron precisamente al escribirlos. Así ocurrió con el
autoguardado de las anotaciones, cuya prueba reveló que una creación en curso podía escribir
su identificador en el examen siguiente (\cref{SEC:IMP-CRITICAL}).

\subsection{Reparto de responsabilidades}

Al autor le correspondieron la definición de los objetivos y del alcance, la elección de las
herramientas y de la forma de trabajo en paralelo, la aceptación o el rechazo de los
resultados y la dirección de producto. El caso más visible fue el rechazo de la primera
dirección visual, que condujo al sistema actual (\cref{SEC:DES-VISUAL}). También le
correspondieron las autorizaciones con efectos externos, como la publicación de la
demostración. A los agentes les correspondió la ejecución: lectura del repositorio y del
contrato, propuesta de soluciones, escritura de código, pruebas y documentación técnica, y
la redacción de borradores de esta memoria. Entre ambos se situó una capa de reglas escritas
(instrucciones de proyecto, mapa de propiedad, guía de estilo y la prohibición de inventar
operaciones de la \acs{api}) que convertía las decisiones del autor en restricciones
comprobables.

\pendiente{El autor debe confirmar y matizar este reparto con su experiencia directa: qué
decisiones tomó personalmente en cada fase, qué propuestas de los agentes rechazó o corrigió
y qué revisó línea a línea.}
